\section{Introduction}

Brain--machine interfaces (BMIs) seek to transform neural activity into useful motor outputs, making continuous trajectory decoding a central problem in prosthetic control and motor neuroscience. In this task, the objective is to estimate the planar hand trajectory of a macaque from spike trains recorded from 98 neural units during reaches to eight target directions. The decoder must remain causal, beginning prediction at 320\,ms and updating every 20\,ms using only current and past observations.

A central challenge is choosing an appropriate level of model structure. Machine learning methods for neural decoding range from simple linear models to highly expressive nonlinear approaches, and method choice remains strongly task-dependent \cite{b1}. At the same time, motor cortical population activity often evolves on low-dimensional manifolds that capture substantial task-relevant structure \cite{b2}. This suggests that improved decoding may depend not only on model complexity, but also on how well the decoder matches the latent structure of the neural data.

Motivated by this trade-off, we first explored a \textit{Hard Direction Decoder}, which performs early direction selection followed by dynamics-based trajectory estimation. We then developed a \textit{Time-Dependent PCA Decoder}, which leverages low-dimensional representations throughout the decoding process and incorporates lightweight causal stabilisation. Our results suggest that, for this dataset, robust performance is achieved more effectively through appropriate low-dimensional structure and pipeline design than through increasing model complexity alone.